\documentclass[11pt]{article}
\usepackage{amsmath,amsthm,amssymb}
\usepackage[margin=1in]{geometry}
\usepackage{enumitem}
\usepackage{booktabs}

\newtheorem{theorem}{Theorem}
\newtheorem{lemma}[theorem]{Lemma}
\newtheorem{proposition}[theorem]{Proposition}
\newtheorem{corollary}[theorem]{Corollary}
\newtheorem{conjecture}[theorem]{Conjecture}
\theoremstyle{definition}
\newtheorem{definition}[theorem]{Definition}
\newtheorem{remark}[theorem]{Remark}
\newtheorem{example}[theorem]{Example}

\title{Structural rank-1 proof for the staircase Bott--Samelson on
$V_{(2,2,1)}$ and $V_{(3,1,1)}$ at $S_5$,\\
plus refutation of the atom-count conjecture}
\author{Clio}
\date{8 May 2026}

\begin{document}
\maketitle

\begin{abstract}
We give a \emph{structural} proof that the staircase Bott--Samelson
operator $\Pi^{S_5}$ has rank exactly $1$ on each of $V_{(2,2,1)}$
and $V_{(3,1,1)}$ at $S_5$. The proof tracks the image dimension
through each Hecke factor $(T_i + 1)$ in the seminormal basis,
exploiting that $T_1$ has a clean eigenvalue spectrum
($\{q, -1\}$ with low multiplicity at the level we need) and that
$(T_1+1)$ kills the $-1$-eigenspace. The argument upgrades the
May~7 ``rank-1 lemma'' (verified by Lagrange interpolation) to a
structural assertion. We extend the rank survey of $\Pi^{S_n}$ to
$n=6$, refute the atom-count conjecture
$\mathrm{rank} \stackrel{?}{=} \lceil k_\lambda/2 \rceil$
(it fails for $\lambda \in \{(4,1,1), (3,3)\}$ at $n=6$), and
record the new data: at $n=6$, ranks $(1,2,3,1,1,2,0,1,0,0,0)$
for the eleven partitions of $6$.
\end{abstract}

\section{Setup}\label{sec:setup}

\subsection{Conventions}

Let $H_q(S_n)$ be the Iwahori--Hecke algebra of type $A_{n-1}$ with
generators $T_1, \ldots, T_{n-1}$ subject to $T_i^2 = (q-1)T_i + q$
and the braid relations. We work over a parameter $q$
(later specialized for verification). For each partition
$\lambda \vdash n$, let $V_\lambda$ denote the irreducible
$H_q(S_n)$-module of shape $\lambda$ (the cell module / Specht module).

\paragraph{Reduced word and the staircase Bott--Samelson.}
The staircase reduced word for the longest element
$w_0 \in S_n$ is
\[
[1, 2, 1,\, 3, 2, 1,\, 4, 3, 2, 1,\, \ldots,\, n{-}1, n{-}2, \ldots, 1].
\]
The corresponding Bott--Samelson element is
\[
\Pi^{S_n} \;:=\; \prod_{k=1}^{n-1}\, (T_k+1)(T_{k-1}+1)\cdots(T_1+1).
\]
At $S_5$,
$\Pi^{S_5} = (T_1{+}1)(T_2{+}1)(T_1{+}1)(T_3{+}1)(T_2{+}1)(T_1{+}1)(T_4{+}1)(T_3{+}1)(T_2{+}1)(T_1{+}1).$
We read products left-to-right: applying $\Pi^{S_5}$ to $v$ means
applying the rightmost factor $(T_1{+}1)$ first.

\paragraph{Seminormal basis.}
For $\lambda \vdash n$, the seminormal basis of $V_\lambda$ is indexed
by standard Young tableaux $T$ of shape $\lambda$, with vectors $v_T$.
For each simple reflection $s_i = (i, i+1)$ and each SYT $T$:
\begin{itemize}[topsep=0.3em,itemsep=0.2em,leftmargin=2em]
\item If $i, i+1$ lie in the same row of $T$ (axial distance $d=1$),
$T_i v_T = q\, v_T$.
\item If $i, i+1$ lie in the same column of $T$ (axial distance $d=-1$),
$T_i v_T = -v_T$.
\item Otherwise (axial distance $|d| \geq 2$), $T_i$ acts as a
$2\times 2$ block on $\mathrm{span}(v_T, v_{s_i T})$ with
eigenvalues $q$ and $-1$. The $+q$-eigenvector is
$\alpha_d\, v_T + \beta_d\, v_{s_i T}$ for specific
coefficients $\alpha_d, \beta_d$ depending only on $d$.
\end{itemize}

The operator $(T_i+1)$ thus acts as:
\begin{itemize}[topsep=0.3em,itemsep=0.2em,leftmargin=2em]
\item Multiplication by $(1+q)$ on $v_T$ if $i, i+1$ in same row.
\item Zero on $v_T$ if $i, i+1$ in same column.
\item A rank-$1$ projector with eigenvalue $(1+q)$ on each $2\times 2$
block (the $+q$-eigenspace direction).
\end{itemize}

\subsection{The key principle for image-tracking}\label{sec:principle}

Composition of $(T_i+1)$ factors yields a sequence of "image-shrinking"
steps. The key principle:
\begin{quote}\itshape
If at some step the current image is contained in
$U \oplus W$ where $U \subseteq +q$-eigenspace of $T_i$ and
$W \subseteq -1$-eigenspace of $T_i$, then applying $(T_i+1)$
collapses to $(1+q)\cdot U$.
\end{quote}
This is the engine of all our rank bounds.

\section{Theorem A: $\mathrm{rank}\,\Pi^{S_5}|_{V_{(2,2,1)}} \leq 1$}
\label{sec:221-rank1}

\subsection{The seminormal basis of $V_{(2,2,1)}$}

The five SYTs of shape $(2,2,1)$ at $n=5$ are:
\begin{align*}
T^{(1)} &= ((1,2),(3,4),(5)), &
T^{(2)} &= ((1,2),(3,5),(4)), \\
T^{(3)} &= ((1,3),(2,4),(5)), &
T^{(4)} &= ((1,3),(2,5),(4)), \\
T^{(5)} &= ((1,4),(2,5),(3)). &
\end{align*}
We write $v_k := v_{T^{(k)}}$.

\subsection{Action of each $T_i$ on the seminormal basis}

Compiling the position data:
\begin{itemize}[topsep=0.3em,itemsep=0.2em,leftmargin=2em]
\item \textbf{$T_1$ (swap $1, 2$).}
  $T_1 v_k = q\, v_k$ for $k=1,2$;
  $T_1 v_k = -v_k$ for $k=3,4,5$.
  Hence $(T_1{+}1)$ has $+q$-eigenspace $\mathrm{span}(v_1, v_2)$
  (dim $2$) and $-1$-eigenspace $\mathrm{span}(v_3, v_4, v_5)$
  (dim $3$).
\item \textbf{$T_2$ (swap $2, 3$).}
  Pairing structure:
  $\{T^{(1)}, T^{(3)}\}$ at axial distance $\pm 2$,
  $\{T^{(2)}, T^{(4)}\}$ at axial distance $\pm 2$,
  and $T^{(5)}$ alone with $T_2 v_5 = -v_5$.
\item \textbf{$T_3$ (swap $3, 4$).}
  $T_3 v_1 = q\, v_1$ (same row, $d=1$);
  $T_3 v_2 = -v_2$ (same column, $d=-1$);
  $T_3 v_3 = -v_3$ ($d=-1$);
  $\{T^{(4)}, T^{(5)}\}$ paired at axial distance $\pm 3$.
\item \textbf{$T_4$ (swap $4, 5$).}
  $\{T^{(1)}, T^{(2)}\}$ paired at $d = \pm 2$;
  $\{T^{(3)}, T^{(4)}\}$ paired at $d = \pm 2$;
  $T_4 v_5 = -v_5$.
\end{itemize}

\subsection{Iterative image computation}

Let $\Pi_k$ denote the operator obtained from the first $k$ factors of
$\Pi^{S_5}$ applied right-to-left. Let $\mathrm{Im}_k := \Pi_k(V_{(2,2,1)})$.

The factors applied right-to-left, corresponding to the staircase word
$[1, 2, 1, 3, 2, 1, 4, 3, 2, 1]$ reversed, are
\[
\underbrace{(T_1{+}1)}_{1}\;
\underbrace{(T_2{+}1)}_{2}\;
\underbrace{(T_3{+}1)}_{3}\;
\underbrace{(T_4{+}1)}_{4}\;
\underbrace{(T_1{+}1)}_{5}\;
\underbrace{(T_2{+}1)}_{6}\;
\underbrace{(T_3{+}1)}_{7}\;
\underbrace{(T_1{+}1)}_{8}\;
\underbrace{(T_2{+}1)}_{9}\;
\underbrace{(T_1{+}1)}_{10}.
\]

\paragraph{Step 1: $(T_1{+}1)$.} Image $\subseteq +q$-eigenspace
of $T_1$, i.e., $\mathrm{Im}_1 \subseteq \mathrm{span}(v_1, v_2)$.
$\dim \mathrm{Im}_1 \leq 2$.

\paragraph{Step 2: $(T_2{+}1)$.}
On the $\{v_1, v_3\}$ block ($d=\pm 2$), $(T_2{+}1) v_1$ is in the
$+q$-eigenspace direction $\alpha_2 v_1 + \beta_2 v_3$
(for Hoefsmit constants $\alpha_2, \beta_2$ depending on $d=2$).
Similarly $(T_2{+}1) v_2 \in \mathrm{span}(\alpha_2 v_2 + \beta_2 v_4)$.
Hence
\[
\mathrm{Im}_2 \;\subseteq\; \mathrm{span}\bigl(\alpha_2 v_1 + \beta_2 v_3,\;
                                    \alpha_2 v_2 + \beta_2 v_4\bigr),
\]
with $\dim \mathrm{Im}_2 \leq 2$.

\paragraph{Step 3: $(T_3{+}1)$.}
$T_3$ acts on the seminormal basis: $v_1$ with eigenvalue $q$ (same row),
$v_2$ with eigenvalue $-1$ (same column), $v_3$ with eigenvalue $-1$
(same column), and $\{v_4, v_5\}$ paired at $d=\pm 3$.
\begin{itemize}[topsep=0.3em,itemsep=0.2em,leftmargin=2em]
\item $(T_3{+}1)(\alpha_2 v_1 + \beta_2 v_3)
= \alpha_2 (1{+}q) v_1 + 0 = \alpha_2(1{+}q) v_1$.
\item $(T_3{+}1)(\alpha_2 v_2 + \beta_2 v_4)
= 0 + \beta_2 (T_3{+}1) v_4
= \beta_2 (1{+}q)\alpha_3 (\alpha_3 v_4 + \beta_3 v_5)$
(rank-1 in $\{v_4, v_5\}$ block of $T_3$).
\end{itemize}
Therefore
\[
\mathrm{Im}_3 \;\subseteq\;
\mathrm{span}\bigl(v_1,\; \alpha_3 v_4 + \beta_3 v_5\bigr),
\quad \dim \mathrm{Im}_3 \leq 2.
\]

\paragraph{Step 4: $(T_4{+}1)$.}
$T_4$ acts: $\{v_1, v_2\}$ paired at $d=\pm 2$;
$\{v_3, v_4\}$ paired at $d=\pm 2$;
$T_4 v_5 = -v_5$.
\begin{itemize}[topsep=0.3em,itemsep=0.2em,leftmargin=2em]
\item $(T_4{+}1) v_1 = (1{+}q)\alpha_2 (\alpha_2 v_1 + \beta_2 v_2)$
(rank-1 in $\{v_1, v_2\}$ block).
\item $(T_4{+}1)(\alpha_3 v_4 + \beta_3 v_5) =
\alpha_3 (T_4{+}1) v_4 + 0
= \alpha_3 (1{+}q)\alpha_2 (\alpha_2 v_3 + \beta_2 v_4)$
(rank-1 in $\{v_3, v_4\}$ block).
\end{itemize}
Therefore
\[
\mathrm{Im}_4 \;\subseteq\;
\mathrm{span}\bigl(\alpha_2 v_1 + \beta_2 v_2,\;
                   \alpha_2 v_3 + \beta_2 v_4\bigr),
\quad \dim \mathrm{Im}_4 \leq 2.
\]

\paragraph{Step 5: $(T_1{+}1)$.}
$T_1$ has $+q$-eigenspace $\mathrm{span}(v_1, v_2)$ and
$-1$-eigenspace $\mathrm{span}(v_3, v_4, v_5)$.
\begin{itemize}[topsep=0.3em,itemsep=0.2em,leftmargin=2em]
\item $(T_1{+}1)(\alpha_2 v_1 + \beta_2 v_2)
= (1{+}q)(\alpha_2 v_1 + \beta_2 v_2)$.
\item $(T_1{+}1)(\alpha_2 v_3 + \beta_2 v_4) = 0 + 0 = 0$.
\end{itemize}
\textbf{The $\mathrm{span}(v_3, v_4)$ component is killed.}
Therefore
\[
\mathrm{Im}_5 \;\subseteq\; (1{+}q)\,\mathrm{span}(\alpha_2 v_1 + \beta_2 v_2),
\quad \dim \mathrm{Im}_5 \leq 1.
\]

\paragraph{Steps 6--10.} Subsequent factors can only preserve or
reduce the image dimension. Therefore
$\dim \mathrm{Im}(\Pi^{S_5}|_{V_{(2,2,1)}}) \leq \dim \mathrm{Im}_5 \leq 1$.

\subsection{The image is nonzero}

By the May~7 verification (Lagrange interpolation through degree-18
$q$-values, recording diagonal entries), the trace
$\sigma_1(V_{(2,2,1)})|_{S_5} = q^4(1+q)^2 \neq 0$. Hence
$\Pi^{S_5}|_{V_{(2,2,1)}}$ is nonzero, so $\dim \mathrm{Im} \geq 1$.

\begin{theorem}[Theorem A: structural rank-1]\label{thm:221}
$\Pi^{S_5}$ acting on $V_{(2,2,1)}$ has rank exactly $1$.
\end{theorem}

\begin{proof}
$\dim \mathrm{Im}_5 \leq 1$ by the iterative image computation
(Steps 1--5). Subsequent factors preserve this bound, so
$\mathrm{rank}\,\Pi^{S_5}|_{V_{(2,2,1)}} \leq 1$. The trace is nonzero,
so $\mathrm{rank} \geq 1$. Hence $\mathrm{rank} = 1$.
\end{proof}

\begin{remark}[Why Step 5 collapses]
The collapse $\dim \mathrm{Im}_4 = 2 \to \dim \mathrm{Im}_5 \leq 1$
exploits two specific features of $V_{(2,2,1)}$:
\begin{itemize}[topsep=0.3em,itemsep=0.2em,leftmargin=2em]
\item After Step 4, the image lies in
$\mathrm{span}(v_1, v_2)\oplus\mathrm{span}(v_3, v_4)$
--- the two pieces sit on opposite sides of the
$T_1$ eigenspace decomposition.
\item $(T_1{+}1)$ at Step 5 then kills $\mathrm{span}(v_3, v_4)$
(both in the $-1$-eigenspace of $T_1$) and preserves
$\mathrm{span}(v_1, v_2)$ (in the $+q$-eigenspace, scaled by $(1{+}q)$).
\end{itemize}
This is the rank-collapsing engine.
\end{remark}

\section{Theorem B: $\mathrm{rank}\,\Pi^{S_5}|_{V_{(3,1,1)}} \leq 1$}
\label{sec:311-rank1}

\subsection{The seminormal basis of $V_{(3,1,1)}$}

The six SYTs of shape $(3,1,1)$ at $n=5$ are:
\begin{align*}
T^{(1)} &= ((1,2,3),(4),(5)), &
T^{(2)} &= ((1,2,4),(3),(5)), &
T^{(3)} &= ((1,2,5),(3),(4)), \\
T^{(4)} &= ((1,3,4),(2),(5)), &
T^{(5)} &= ((1,3,5),(2),(4)), &
T^{(6)} &= ((1,4,5),(2),(3)).
\end{align*}
We write $v_k := v_{T^{(k)}}$.

\subsection{Action of each $T_i$}

Position lookups give:
\begin{itemize}[topsep=0.3em,itemsep=0.2em,leftmargin=2em]
\item \textbf{$T_1$.} $T_1 v_k = q\, v_k$ for $k=1,2,3$;
  $T_1 v_k = -v_k$ for $k=4,5,6$.
  Hence $+q$-eigenspace of $T_1$ is $\mathrm{span}(v_1, v_2, v_3)$ (dim $3$).
\item \textbf{$T_2$.}
  $T_2 v_1 = q\, v_1$;
  $\{T^{(2)}, T^{(4)}\}$ paired at $d=\pm 2$;
  $\{T^{(3)}, T^{(5)}\}$ paired at $d=\pm 2$;
  $T_2 v_6 = -v_6$.
\item \textbf{$T_3$.}
  $\{T^{(1)}, T^{(2)}\}$ paired at $d=\pm 3$;
  $T_3 v_3 = -v_3$;
  $T_3 v_4 = q\, v_4$;
  $\{T^{(5)}, T^{(6)}\}$ paired at $d=\pm 3$.
\item \textbf{$T_4$.}
  $T_4 v_1 = -v_1$;
  $\{T^{(2)}, T^{(3)}\}$ paired at $d=\pm 4$;
  $\{T^{(4)}, T^{(5)}\}$ paired at $d=\pm 4$;
  $T_4 v_6 = q\, v_6$.
\end{itemize}

\subsection{Iterative image computation}

We apply factors of $\Pi^{S_5}$ from right to left: the order is
\[
\underbrace{(T_1{+}1)}_{\text{Step 1}}\;
\underbrace{(T_2{+}1)}_{\text{Step 2}}\;
\underbrace{(T_3{+}1)}_{\text{Step 3}}\;
\underbrace{(T_4{+}1)}_{\text{Step 4}}\;
\underbrace{(T_1{+}1)}_{\text{Step 5}}\;
\underbrace{(T_2{+}1)\cdots(T_1{+}1)}_{\text{Steps 6--10}}.
\]

The proof tracks $\mathrm{Im}_k$ very carefully through Steps~1--5,
showing $\dim \mathrm{Im}_5 \leq 1$. The remaining factors then preserve
this bound.

\paragraph{Step 1: $(T_1{+}1)$.}
$\mathrm{Im}_1 \subseteq \mathrm{span}(v_1, v_2, v_3)$; $\dim \leq 3$.

\paragraph{Step 2: $(T_2{+}1)$.}
$(T_2{+}1) v_1 = (1{+}q) v_1$ (same row).
$(T_2{+}1) v_2 \in (1{+}q)\,\mathrm{span}(\alpha_2 v_2 + \beta_2 v_4)$
(rank-1 $+q$-direction in $d{=}\pm 2$ block).
$(T_2{+}1) v_3 \in (1{+}q)\,\mathrm{span}(\alpha_2 v_3 + \beta_2 v_5)$.
Hence $\mathrm{Im}_2$ is generated by the three vectors
\[
b_1 = v_1, \quad
b_2 = \alpha_2 v_2 + \beta_2 v_4, \quad
b_3 = \alpha_2 v_3 + \beta_2 v_5,
\]
each up to a $(1{+}q)$ scalar.

\paragraph{Step 3: $(T_3{+}1)$.}
\begin{itemize}[topsep=0.3em,itemsep=0.2em,leftmargin=2em]
\item $(T_3{+}1) b_1 = (T_3{+}1) v_1 = (1{+}q)\alpha_3
(\alpha_3 v_1 + \beta_3 v_2)$
(rank-1 in the $\{v_1, v_2\}$ block of $T_3$, axial $d=\pm 3$).
\item $(T_3{+}1) b_2 = \alpha_2 (T_3{+}1) v_2 + \beta_2 (1{+}q) v_4$.
Now $(T_3{+}1) v_2 \in (1{+}q)\,\mathrm{span}(\alpha_3 v_1 + \beta_3 v_2)$
(same block; $v_2$ has $d{=}+3$). So
$(T_3{+}1) b_2 \in
\mathrm{span}(\alpha_3 v_1 + \beta_3 v_2)\;\oplus\;\mathrm{span}(v_4)$.
\item $(T_3{+}1) b_3 = 0 + \beta_2 (T_3{+}1) v_5$.
$(T_3{+}1) v_5 \in (1{+}q)\,\mathrm{span}(\alpha_3 v_5 + \beta_3 v_6)$
(rank-1 in $\{v_5, v_6\}$ block).
\end{itemize}
Hence $\mathrm{Im}_3$ is generated (up to scalars) by three vectors
\[
a_1 = \alpha_3 v_1 + \beta_3 v_2,\quad
a_2 \in \mathrm{span}(\alpha_3 v_1 + \beta_3 v_2)\oplus\mathrm{span}(v_4),\quad
a_3 = \alpha_3 v_5 + \beta_3 v_6.
\]
The key structural fact: the projection of $\mathrm{Im}_3$ onto
$\mathrm{span}(v_1, v_2, v_3)$ is contained in
$\mathrm{span}(\alpha_3 v_1 + \beta_3 v_2)$ --- a $1$-dim line, NOT
the full $\mathrm{span}(v_1, v_2)$.

\paragraph{Step 4: $(T_4{+}1)$.}
Recall $T_4$ acts:
$T_4 v_1 = -v_1$;
rank-1 projection on $\{v_2, v_3\}$ at $d=\pm 4$;
rank-1 projection on $\{v_4, v_5\}$ at $d=\pm 4$;
$T_4 v_6 = q v_6$.

\begin{itemize}[topsep=0.3em,itemsep=0.2em,leftmargin=2em]
\item $(T_4{+}1) a_1 = \alpha_3 (T_4{+}1) v_1 + \beta_3 (T_4{+}1) v_2
                      = 0 + \beta_3 (1{+}q)\alpha_4(\alpha_4 v_2 + \beta_4 v_3)$.
So $(T_4{+}1) a_1 \in \mathrm{span}(\alpha_4 v_2 + \beta_4 v_3)$
\quad\textbf{(1-dim)}.
\item $(T_4{+}1) a_2$: since $a_2 \in \mathrm{span}(\alpha_3 v_1
+ \beta_3 v_2)\oplus\mathrm{span}(v_4)$, and $(T_4{+}1) v_1 = 0$,
$(T_4{+}1) v_2 \in \mathrm{span}(\alpha_4 v_2 + \beta_4 v_3)$,
$(T_4{+}1) v_4 \in \mathrm{span}(\alpha_4 v_4 + \beta_4 v_5)$:
$(T_4{+}1) a_2 \in \mathrm{span}(\alpha_4 v_2 + \beta_4 v_3,\;
                                 \alpha_4 v_4 + \beta_4 v_5)$
\quad\textbf{(2-dim)}.
\item $(T_4{+}1) a_3$: $a_3 \in \mathrm{span}(v_5, v_6)$,
$(T_4{+}1) v_5 \in \mathrm{span}(\alpha_4 v_4 + \beta_4 v_5)$,
$(T_4{+}1) v_6 = (1{+}q) v_6$.
$(T_4{+}1) a_3 \in \mathrm{span}(\alpha_4 v_4 + \beta_4 v_5,\;v_6)$
\quad\textbf{(2-dim)}.
\end{itemize}

Therefore $\mathrm{Im}_4 \subseteq \mathrm{span}(\alpha_4 v_2 + \beta_4 v_3,\;
\alpha_4 v_4 + \beta_4 v_5,\;v_6)$.
Crucially, the projection of $\mathrm{Im}_4$ onto
$\mathrm{span}(v_1, v_2, v_3)$ lies in
$\mathrm{span}(\alpha_4 v_2 + \beta_4 v_3)$ \textbf{(1-dim)}.

\paragraph{Step 5: $(T_1{+}1)$.}
$T_1$ kills $\mathrm{span}(v_4, v_5, v_6)$ (the $-1$-eigenspace) and
acts as $(1{+}q)$ on $\mathrm{span}(v_1, v_2, v_3)$. Therefore
\[
(T_1{+}1)\,\mathrm{Im}_4 \;\subseteq\; (1{+}q)\,
\mathrm{span}\bigl(\alpha_4 v_2 + \beta_4 v_3\bigr),
\]
\textbf{a $1$-dim subspace}.

\paragraph{Steps 6--10.} Each $(T_i{+}1)$ factor is a linear map and
cannot increase image dimension. So
$\dim \mathrm{Im}(\Pi^{S_5}|_{V_{(3,1,1)}}) \leq \dim \mathrm{Im}_5 \leq 1$.

\begin{theorem}[Theorem B: structural rank-1]\label{thm:311}
$\Pi^{S_5}$ acting on $V_{(3,1,1)}$ has rank exactly $1$.
\end{theorem}

\begin{proof}
$\dim \mathrm{Im}_5 \leq 1$ by the refined argument above:
the $v_4, v_5, v_6$ components of $\mathrm{Im}_4$ are killed by
$(T_1{+}1)$, leaving only the $\mathrm{span}(\alpha_4 v_2 + \beta_4 v_3)$
direction, which is $1$-dimensional. Subsequent factors preserve
this bound.

The image is nonzero by the May~7 verification:
$\sigma_1(V_{(3,1,1)})|_{S_5} = q^3(1{+}q)^4 + 2q^4(1{+}q)^2 \neq 0$.
Hence $\mathrm{rank} = 1$.
\end{proof}

\section{The structural argument vs.\ rank-2 cases}\label{sec:rank2}

For comparison, we apply the same image-tracking to
$V_{(3,2)}$ at $S_5$ (rank $2$).

The five SYTs of $(3,2)$ are
$T^{(1)} = ((1,2,3),(4,5))$,
$T^{(2)} = ((1,2,4),(3,5))$,
$T^{(3)} = ((1,2,5),(3,4))$,
$T^{(4)} = ((1,3,4),(2,5))$,
$T^{(5)} = ((1,3,5),(2,4))$.

In particular, $T_3$ acts on $V_{(3,2)}$ with the structure:
\begin{itemize}[topsep=0.3em,itemsep=0.2em,leftmargin=2em]
\item $\{v_1, v_2\}$ paired at $d=\pm 3$.
\item $T_3 v_3 = q\, v_3$ (since $3, 4$ in same row of $T^{(3)}$).
\item $T_3 v_4 = q\, v_4$ (same row).
\item $T_3 v_5 = -v_5$ (same column).
\end{itemize}

\paragraph{Crucial difference from $V_{(3,1,1)}$.}
In $V_{(3,1,1)}$, $T_3 v_3 = -v_3$ (since the
shape has $3, 4$ in the same column for $T^{(3)}$). Therefore
in Step 3 of the $V_{(3,1,1)}$ argument, the
$\alpha_2 v_3 + \beta_2 v_5$ term loses its $v_3$-component:
\[
(T_3{+}1)(\alpha_2 v_3 + \beta_2 v_5)\big|_{V_{(3,1,1)}}
\;=\; 0 + \beta_2 (T_3{+}1) v_5
\;\in\; \mathrm{span}(v_5, v_6),
\]
and this entire contribution sits in the $-1$-eigenspace of $T_1$,
so it gets killed by Step 5.

In $V_{(3,2)}$, by contrast, $T_3 v_3 = q v_3$, so
\[
(T_3{+}1)(\alpha_2 v_3 + \beta_2 v_5)\big|_{V_{(3,2)}}
\;=\; \alpha_2 (1{+}q) v_3 + 0
\;=\; \alpha_2(1{+}q)\, v_3.
\]
This contribution stays in $\mathrm{span}(v_3) \subseteq +q\text{-eig}(T_1)$,
\textbf{surviving} the Step 5 application of $(T_1{+}1)$.

Tracking through all $10$ steps numerically (script
\texttt{verify\_seminormal.py}) shows that for $V_{(3,2)}$
the image dimension is $3$ through Step $4$, then drops to $2$ at Step $5$
and stays at $2$. The rank is $2$, not $1$.

\begin{remark}[Where rank-1 vs rank-2 is decided]
The image dimension sequence at Steps $4 \to 5$ encodes the rank:
\begin{itemize}[topsep=0.3em,itemsep=0.2em,leftmargin=2em]
\item $V_{(2,2,1)}$: $2 \to 1$. Image after Step 4 has $\mathrm{span}(v_1, v_2)$
component $1$-dim; the rest is killed by $(T_1{+}1)$.
\item $V_{(3,1,1)}$: $3 \to 1$. Image after Step 4 has $+q$-eig$(T_1)$
component $1$-dim (a single line $\alpha_4 v_2 + \beta_4 v_3$); the rest
is in $-1$-eig$(T_1)$ and killed.
\item $V_{(3,2)}$: $3 \to 2$. Image after Step 4 has $+q$-eig$(T_1)$
component $2$-dim; nothing more collapses at Step 5.
\item $V_{(4,1)}$: similarly $3 \to 2$.
\end{itemize}
\end{remark}

\section{Extended rank survey: $n = 6$}\label{sec:n6}

We extend the rank survey of \texttt{rank\_survey.py} (May~7) to $n = 6$.
The script \texttt{rank\_n6\_survey.py} computes
$\Pi^{S_6}$ on each cell module $V_\lambda$ via Lagrange interpolation
and reports the numerical rank at $q \in \{2,3,5,7,11,13\}$.

\subsection{Data}

\begin{center}
\begin{tabular}{l c c c}
\toprule
$\lambda \vdash 6$ & $\dim V_\lambda$ & $\mathrm{rank}\,\Pi^{S_6}$ & $\sigma_1(V_\lambda)$ (degrees only) \\
\midrule
$(6)$         &  1 & 1 & $1$\,..\,$q^{15}$ \\
$(5, 1)$      &  5 & 2 & $q$\,..\,$q^{14}$ \\
$(4, 2)$      &  9 & 3 & $q^2$\,..\,$q^{13}$ \\
$(4, 1, 1)$   & 10 & 1 & $q^3$\,..\,$q^{12}$ \\
$(3, 3)$      &  5 & 1 & $q^3$\,..\,$q^{12}$ \\
$(3, 2, 1)$   & 16 & 2 & $q^4$\,..\,$q^{11}$ \\
$(3, 1, 1, 1)$& 10 & 0 & $0$ \\
$(2, 2, 2)$   &  5 & 1 & $q^6$\,..\,$q^9$ \\
$(2, 2, 1, 1)$&  9 & 0 & $0$ \\
$(2, 1, 1, 1, 1)$ & 5 & 0 & $0$ \\
$(1^6)$       &  1 & 0 & $0$ \\
\bottomrule
\end{tabular}
\end{center}

\subsection{Coincidence: $\sigma_1((4,1,1)) = \sigma_1((3,3))$}

A striking observation: the trace polynomials are equal:
\[
\sigma_1(V_{(4,1,1)}) \;=\; \sigma_1(V_{(3,3)}) \;=\;
q^3 + 15 q^4 + 91 q^5 + 281 q^6 + 484 q^7 + 484 q^8 + 281 q^9 + 91 q^{10} + 15 q^{11} + q^{12}.
\]
Both partitions have rank $1$, but the dimensions $10$ and $5$ are different.
This is unique among $n \leq 6$: at $n=5$, all $\sigma_1$ values are
distinct.

\subsection{Atom decomposition}

Decomposing each nonzero $\sigma_1$ as
$\sum_{i \geq 0} m^{(i)}\, q^{n(\lambda) + i}\,(1+q)^{w_\lambda - 2i}$
where $w_\lambda = L - 2 n(\lambda)$ is the width:

\begin{center}
\begin{tabular}{l c c l}
\toprule
$\lambda$ & rank & $k_\lambda$ & atom multiplicities $(m^{(0)}, m^{(1)}, \ldots)$ \\
\midrule
$(6)$         & 1 & 1 & $(1)$ \\
$(5, 1)$      & 2 & 4 & $(1, 3, 9, 14, 0, 0, 0)$ \\
$(4, 2)$      & 3 & 5 & $(1, 5, 21, 30, 9, 0)$ \\
$(4, 1, 1)$   & 1 & 4 & $(1, 6, 13, 6, 0)$ \\
$(3, 3)$      & 1 & 4 & $(1, 6, 13, 6, 0)$ \\
$(3, 2, 1)$   & 2 & 4 & $(1, 9, 15, 2)$ \\
$(2, 2, 2)$   & 1 & 1 & $(1, 0)$ \\
\bottomrule
\end{tabular}
\end{center}

(Here $k_\lambda$ denotes the number of nonzero atoms.)

\section{Refutation of the May-7 atom-count conjecture}\label{sec:refute}

\begin{conjecture}[\emph{Conjecture from May-7 writeup, hereby refuted}]
\label{conj:atom-half}
$\mathrm{rank}\,\Pi^{S_n}|_{V_\lambda}\;=\;\lceil k_\lambda / 2 \rceil$.
\end{conjecture}

\begin{proposition}[Refutation]
At $n = 6$, $V_{(4,1,1)}$ and $V_{(3,3)}$ each have $k_\lambda = 4$
nonzero atoms but $\mathrm{rank}\,\Pi^{S_6} = 1 \neq 2 = \lceil 4/2 \rceil$.
\end{proposition}

\begin{proof}
By the rank survey of Section~\ref{sec:n6} (computational), and the
atom decomposition in the same section.
\end{proof}

\paragraph{What survives.} The conjecture holds at $n \leq 5$
(by direct check) and at $n=6$ for $\lambda \in \{(6), (5,1), (4,2),
(3,2,1), (2,2,2)\}$. It fails specifically for the two coinciding
$\sigma_1$ partitions $\{(4,1,1), (3,3)\}$.

\paragraph{Refined conjecture.}
\begin{conjecture}[Atom count is an upper bound]\label{conj:atom-bound}
$\mathrm{rank}\,\Pi^{S_n}|_{V_\lambda}\;\leq\;\lceil k_\lambda / 2 \rceil$.
\end{conjecture}
This holds for all $n \leq 6$ data points and matches the original
spirit of the May-7 conjecture as a one-sided bound.

\section{Discussion and open questions}\label{sec:open}

\subsection{What the structural argument captures}

The iterative-image-tracking argument of
Sections~\ref{sec:221-rank1}--\ref{sec:311-rank1}, which gives
rank-$1$ structurally for $V_{(2,2,1)}$ and $V_{(3,1,1)}$, exploits
two key features of these specific cell modules:
\begin{enumerate}[leftmargin=2em,topsep=0.3em,itemsep=0.2em]
\item \emph{Eigenvalue uniformity of $T_1$.}
$T_1$ is diagonal with at most two distinct eigenvalues
($\{q, -1\}$). The $-1$-eigenspace is $\mathrm{span}(v_T : 1 \in D_L(T))$
where $D_L$ is the descent set of the SYT.
\item \emph{The ``narrow'' structure of $\mathrm{Im}_4$ in the $+q$
direction of $T_1$.} For these partitions, $\mathrm{Im}_4$ has
its $+q$-component in a specific 1-dim subspace of
$\mathrm{span}(v_2, v_3)$ (case $V_{(3,1,1)}$) or in
$\mathrm{span}(v_1, v_2)$ with the $v_2$ direction killed by $T_3$
(case $V_{(2,2,1)}$).
\end{enumerate}

The structural argument \emph{fails} for $V_{(3,2)}$ and $V_{(4,1)}$
because $T_3$ acts uniformly on more of the surviving generators of
$\mathrm{Im}_4$, preventing the rank-1 collapse.

\subsection{Open: a structural characterization}

The data through $n=6$ does not yet determine a clean
representation-theoretic invariant equal to
$\mathrm{rank}\,\Pi^{S_n}|_{V_\lambda}$. Possible candidates:
\begin{itemize}[leftmargin=2em,topsep=0.3em,itemsep=0.2em]
\item \emph{Number of ``good'' SYTs}: SYTs $T$ such that
$T_1, T_2, T_3, \ldots$ act ``compatibly'' on $v_T$ in some
specific sense yet to be made precise. Empirically, for each
rank-$r$ case, the first $r$ SYTs in the cell-module enumeration
order give linearly independent images under $\Pi$.
\item \emph{Sub-family rank formulas.}
For $\lambda = (n-2, 2)$: $\mathrm{rank} = n - 3$ (verified for
$n \in \{4,5,6\}$). For $\lambda = (n-1, 1)$: $\mathrm{rank} =
\lceil (n-2)/2 \rceil$ (verified for $n \in \{3,4,5,6\}$).
These are not in general $\lceil k_\lambda/2 \rceil$ but match the
rank exactly for these families.
\end{itemize}

\subsection{Implications for the May-7 trace identity}

Theorems~\ref{thm:221} and~\ref{thm:311} upgrade the May-7
``rank-1 lemma'' (Theorem~3.1 in
\texttt{2026-05-07-rank-one-Pi.tex}) from a computational verification
to a structural theorem. By Corollary~3.7 of that writeup, this
implies the branching factorization
\[
T_A \;=\; \mathrm{tr}_{V_{(3,2,1)}}\!\bigl(\Pi^{S_5}\,P_{(2,2,1)}^{\mathrm{KL}}\,R_5'\bigr) \;=\; q^2\,\sigma_1(V_{(2,2,1)})\big|_{S_5}
\]
and the analogous identity for $T_B$ are now established at full rigor
(modulo the standard branching reduction, which is purely
representation-theoretic).

\section{Honest assessment}\label{sec:honest}

\paragraph{What this paper achieves.}
\begin{enumerate}[leftmargin=2em,topsep=0.3em,itemsep=0.2em]
\item Structural proofs (Theorems~\ref{thm:221}, \ref{thm:311})
that $\Pi^{S_5}$ has rank exactly $1$ on $V_{(2,2,1)}$ and $V_{(3,1,1)}$.
The proofs track the image dimension through each Hecke factor
$(T_i+1)$ in the seminormal basis.
\item Extension of the rank survey to $n=6$
(Section~\ref{sec:n6}), with the new data point that
$\sigma_1(V_{(4,1,1)}) = \sigma_1(V_{(3,3)})$ at $n=6$.
\item Refutation of the May-7 atom-count conjecture
(Section~\ref{sec:refute}), with a refined upper-bound
conjecture in place.
\end{enumerate}

\paragraph{What this paper does \emph{not} achieve.}
\begin{enumerate}[leftmargin=2em,topsep=0.3em,itemsep=0.2em]
\item \emph{A general structural characterization} of
$\mathrm{rank}\,\Pi^{S_n}|_{V_\lambda}$ for all $(n, \lambda)$.
The structural argument used here is special to the small
cases where $T_3$ or $T_4$ has a clean ``splitting'' eigenvalue
behavior on the surviving image generators.
\item \emph{Why the rank survey at $n=6$ has the specific values it does.}
For instance, why does $V_{(4,2)}$ have rank exactly $3$ (and not $2$
or $4$)? The structural argument gives only an upper bound consistent
with rank $3$, but does not pin it down structurally.
\item \emph{An explicit closed form for the rank-$1$ image vector.}
The image of $\Pi^{S_5}$ on $V_{(2,2,1)}$ is spanned by a specific
combination of $v_1, v_2$ with coefficients determined by the
$d=\pm 2$ Hoefsmit constants; we have not written this combination
explicitly in this writeup.
\end{enumerate}

\paragraph{Status.} The branching factorization conjecture is now a
\emph{structural theorem} (rather than ``theorem-modulo-rank-1-lemma''
as it was in the May-7 writeup). The general rank-characterization
problem remains open, with a partial classification at $n \leq 6$
and refined upper-bound conjectures.

\paragraph{Computational note.} All matrix entries, rank verifications,
and the new $n = 6$ data are checked by Lagrange interpolation
in the scripts:
\begin{center}
\texttt{\textasciitilde/projects/scratch/2026-05-08-rank-Pi/rank\_n6\_survey.py},\\
\texttt{\textasciitilde/projects/scratch/2026-05-08-rank-Pi/image\_analysis.py}.
\end{center}

\end{document}
